%------------------------------------------------------------------------------%
\documentclass[a4paper,12pt,fleqn]{article}

\usepackage{calculo_varias_variaveis-1}

\ShowAnswers

%------------------------------------------------------------------------------%
\begin{document}

\makeHeader{CFVVI}{Exame -- Turma 2}{12/01/2026}

% 01
%------------------------------------------------------------------------------%
\exercise{40\%}
Determine os pontos candidatos a extremos de
\(
  x^2+e^{y}
\)
sujeita à restrição
\(
  x^2+y=1
\)
%\clearpagequestiononly

\begin{answer}
\begin{multicols}{2}
Queremos os extremos de
\[
  f(x, y) = x^2 + e^{y}
\]
sujeita à restrição
\[
  g(x, y) = x^2 + y - 1 = 0
\]
Gradientes
\begin{align*}
  \nabla f & = \vetor{2x}{e^{y}} \\
  \nabla g & = \vetor{2x}{1}
\end{align*}

O sistema de Lagrange é
\begin{align*}
  2x    & =2\lambda x \\
  e^{y} & =\lambda    \\
  x^2+y & =1
\end{align*}

Da primeira equação
\[
  x = 0 \quad \text{ou} \quad \lambda=1
\]

Caso $x=0$, da restrição
\[
  y = 1
\]
e da segunda equação
\[
  \lambda = e
\]
obtermos o primeiro ponto
\[
  P_1 = (0, 1)
\]

Caso $\lambda=1$, da segunda equação
\[
  e^{y}=1 \quad\Rightarrow\quad y=0
\]
Da restrição
\[
x^2=1 \quad\Rightarrow\quad x=\pm1
\]
Encontramos o segundo e terceiro pontos
\[
P_2 = (1, 0) \qquad P_3 = (-1, 0)
\]
\end{multicols}
\end{answer}

% 02
%------------------------------------------------------------------------------%
\exercise{20\%}
Seja \(f(x,y)=e^{x}\ln(1+x^{2}+y^{2})\). Determine a
aproximação linear de \(f\) no ponto \((0, 1)\) e use-a para aproximar
\(f(0.05, 0.98)\).


\begin{answer}
A aproximação linear é
\[
  L(x, y) = f(a, b) + f_x(a, b) (x-a) + f_y(a, b) (y-b)
\]
Cálculo das derivadas parciais
\begin{align*}
  f_x(x, y)
  & = \D{}{x} \left[e^{x}\ln(1+x^{2}+y^{2})\right] \\
  & = \D{}{x} \left[e^{x}\right]\ln(1+x^{2}+y^{2})
    + e^{x}\D{}{x} \left[\ln(1+x^{2}+y^{2})\right] \\
  & = e^{x}\ln(1+x^{2}+y^{2})
    + e^{x} \frac{1}{1+x^{2}+y^{2}} \D{}{x} \left(1+x^{2}+y^{2}\right)\\
  & = e^x\left(\ln\left(1+x^2+y^2\right) + \frac{2x}{1+x^2+y^2}\right)
  \\[3mm]
  f_y(x,y)
  & = \D{}{y} \left[e^{x}\ln(1+x^{2}+y^{2})\right] \\
  & = e^{x}\D{}{y} \left[\ln(1+x^{2}+y^{2})\right] \\
  & = e^{x} \frac{1}{1+x^{2}+y^{2}} \D{}{y} \left(1+x^{2}+y^{2}\right)\\
  & = \frac{2ye^x}{1+x^2+y^2}
\end{align*}
Avaliação em \((0, 1)\)
\begin{align*}
  f(0, 1)
  & = \left[e^{x}\ln(1+x^{2}+y^{2})\right] \evalat{(0, 1)}{} \\
  & = e^{x}\ln(1+0^{2}+1^{2}) \\
  & = \ln 2 \\
  f_x(0, 1)
  & = \left[e^x\left(\ln\left(1+x^2+y^2\right) + \frac{2x}{1+x^2+y^2}\right)\right] \evalat{(0, 1)}{} \\
  & = e^0\left(\ln\left(1+0^2+1^2\right) + \frac{2\times 0}{1+0^2+1^2}\right) \\
  & = \ln 2 \\
  f_y(0, 1)
  & = \left[\frac{2ye^x}{1+x^2+y^2}\right] \evalat{(0, 1)}{} \\
  & = \frac{2\times 1 \times e^0}{1+0^2+1^2} \\
  & = 1
\end{align*}
A aproximação linear é
\begin{align*}
  L(x, y)
  & = \ln(2)    + \ln(2)\,x + (y-1) \\
  & = \ln(2)\,x + y + (\ln(2) -1)
\end{align*}
Aproximação em \((0.05,\,0.98)\)
\begin{align*}
  L(0.05, 0.98)
  & = \ln(2) + \ln(2)\,0.05 + (0.98-1) \\
  & = \ln(2) + \ln(2)\,0.05 - 0.02 \\
  & = 1.05\ln(2) - 0.02
\end{align*}
Com uma calculadora podemos avaliar que
\begin{align*}
  L(0.05, 0.98) & = 0.7078 \\
  f(0.05, 0.98) & = 0.7090
\end{align*}
\end{answer}


% 03
%------------------------------------------------------------------------------%
\exercise{40\%}
Determine e classifique os pontos críticos da função
\(
  f(x,y) = e^{xy} - x^2 - y^2
\)

\clearpagequestiononly

\begin{answer}
  Calculando o gradiente
  \[
    f_x(x, y)
    = y e^{xy} - 2x
  \]
  \[
    f_y(x, y)
    = x e^{xy} - 2y
  \]
  Igualando o gradiente ao vetor nulo
  \[
    \begin{cases}
      y e^{xy} - 2x = 0\\
      x e^{xy} - 2y = 0
    \end{cases}
  \]
  rearranjando
  \[
    \begin{cases}
      2x = y e^{xy} \\
      2y = x e^{xy}
    \end{cases}
  \]
  Subtraindo as equações
  \[
    2(x-y) = -(x-y)e^{xy}
  \]

  Caso 1: $x\neq y$, dividimos os dois lados por $x-y$, obtendo
  \(
    2 = -e^{xy}
  \)
  que não possui solução.

  Caso 2: $x=y$, a equação \(2(x-y) = -(x-y)e^{xy}\) é satisfeita

  Substituindo na primeira equação e obtemos
  \begin{align*}
    2x & = y e^{xy} \\
    2x & = x e^{x^2}
  \end{align*}

  Se $x=0$, essa equação é satisfeita e temos o primeiro ponto crítico \(P_1=(0, 0)\)

  Se $x\neq y$, dividimos os dois termos por $x$
  \begin{align*}
    2x      & = x e^{x^2}        \\
    2       & = e^{x^2}          \\
    \ln(2)  & = x^2              \\
    \abs{x} & = \sqrt{\ln(2)}    \\
    x       & = \pm\sqrt{\ln(2)}
  \end{align*}
  Encontramos os pontos críticos
  \(P_2 = \left( \sqrt{\ln(2)},  \sqrt{\ln(2)}\right)\) e
  \(P_3 = \left(-\sqrt{\ln(2)}, -\sqrt{\ln(2)}\right)\)

  \bigskip
  Calculando a Hessiana de $f$
  \[
    f_{xx}(x, y)
    = y^2 e^{xy}-2
  \]
  \[
    f_{yy}(x, y)
    = x^2 e^{xy}-2
  \]
  \[
    f_{xy}(x, y)
    = xy e^{xy}+e^{xy}
  \]
  Discriminante (Determinante da Hessiana)
  \[
    D(x, y) = f_{xx}(x, y)f_{yy}(x, y) - f^2_{xy}(x, y)
  \]

  \bigskip
  Classificando o ponto \((0, 0)\)
  \[
    f_{xx}(0, 0) = -2
  \]
  \[
    f_{yy}(0, 0) = -2
  \]
  \[
    f_{xy}(0, 0) =  1
  \]
  \[
    D(0, 0) = 4-1 = 3 > 0
  \]
  Como \(f_{xx}(0, 0) = -2 < 0\), \textbf{o ponto \((0,0)\) é um máximo local}

  \bigskip
  No ponto \(\left(\sqrt{\ln 2}, \sqrt{\ln 2}\right)\)
  \[
    f_{xx}\left(\sqrt{\ln 2}, \sqrt{\ln 2}\right)
    = 2\ln 2 -2
  \]
  \[
    f_{yy}\left(\sqrt{\ln 2}, \sqrt{\ln 2}\right)
    = 2\ln 2 -2
  \]
  \[
    f_{xy}\left(\sqrt{\ln 2}, \sqrt{\ln 2}\right)
    = 2\ln 2 +2
  \]
  \[
    D\left(\sqrt{\ln 2}, \sqrt{\ln 2}\right)
    = (2\ln 2 -2)^2 - (2\ln 2 +2)^2
    < 0
  \]
  Portanto, é um \textbf{ponto de sela}

  \bigskip
  No ponto \(\left(-\sqrt{\ln 2}, -\sqrt{\ln 2}\right)\)
  \[
    f_{xx}\left(-\sqrt{\ln 2}, -\sqrt{\ln 2}\right)
    = 2\ln 2 -2
  \]
  \[
    f_{yy}\left(-\sqrt{\ln 2}, -\sqrt{\ln 2}\right)
    = 2\ln 2 -2
  \]
  \[
    f_{xy}\left(-\sqrt{\ln 2}, -\sqrt{\ln 2}\right)
    = 2\ln 2 +2
  \]
  \[
    D\left(-\sqrt{\ln 2}, -\sqrt{\ln 2}\right)
    = (2\ln 2 -2)^2 - (2\ln 2 +2)^2
    < 0
  \]
  Portanto, é um \textbf{ponto de sela}
\end{answer}


%------------------------------------------------------------------------------%
\end{document}
%------------------------------------------------------------------------------%
